\documentclass{article}
\usepackage[utf8]{inputenc}
\usepackage{geometry}
\geometry{
  a4paper,
  total={170mm,257mm},
  left=20mm,
  top=20mm,
}
\usepackage{graphicx}
\usepackage{titling}
\usepackage{booktabs}
\usepackage{listings}
\usepackage{xcolor}
\usepackage{hyperref}
\usepackage{fancyvrb}
\usepackage{float}
\usepackage{amsmath}

\definecolor{darkblue}{rgb}{0.0, 0.0, 0.5}
\RecustomVerbatimEnvironment{verbatim}{Verbatim}{formatcom=\color{darkblue}}
\setlength{\parindent}{0pt}

\title{\textbf{DEEP LEARNING} \\
{\Large \textbf{Report 1: Gradient Descent}}}
\date{May 2026}

\usepackage{fancyhdr}
\fancypagestyle{plain}{
  \fancyhf{}
  \fancyfoot[L]{\thedate}
  \fancyfoot[R]{\thepage}
  \fancyhead[L]{University of Science and Technology of Hanoi}
  \fancyhead[R]{Deep Learning}
}
\makeatletter
\def\@maketitle{%
  \newpage
  \vskip 1em
  \begin{center}
  \let \footnote \thanks
    {\LARGE \@title \par}%
    \vskip 1em
  \end{center}
  \par
  \vskip 1em}
\makeatother

\begin{document}
\pagestyle{plain}
\maketitle

\noindent\begin{tabular}{@{}ll}
  Student: & Vu Ha Vy -- 2540112 \\
\end{tabular}

\section*{Implementation}

The gradient descent function is implemented using a \texttt{while} loop with a convergence threshold:

\begin{verbatim}
def gradient_descent(x0, lr, max_iter=1000, threshold=0.01):
    x = x0
    t = 1
    while t <= max_iter:
        fx = f(x)
        print(f"Time = {t} | x = {x:.4f} | f(x) = {fx:.4f}")
        x_new = x - lr * df(x)
        fx_new = f(x_new)
        if abs(fx_new - fx) < threshold:
            print(f"\n Converged after {t} iterations, f(x) = {fx_new:.4f}")
            break
        x = x_new
        t += 1
    else:
        print(f"\n Reached max iterations ({max_iter}), f(x) = {f(x):.4f}")
\end{verbatim}

The experiment uses $x_0 = 10$, \texttt{max\_iter} = 1000, \texttt{threshold} = 0.01, with four different learning rates
for the function $f(x) = x^2$ ($df(x) = 2x$).

\section*{Analyze of Different Learning Rates}

\textbf{Large learning rate ($lr = 1.0$)}\\
With $lr = 1$, $x$ flips sign every step between 10 and -10. Since $f(x) = x^2$ gives the same value for $+x$ and $-x$, the change in $f(x)$ is exactly zero and the threshold is met after just 1 step, but at $f(x) = 100$ is definitely not the minimum value.
 
\begin{verbatim}
Time = 1 | x = 10.0000 | f(x) = 100.0000

 Converged after 1 iterations, f(x) = 100.0000
\end{verbatim}
\vspace{0.5em}

\textbf{Good learning rate ($lr = 0.1$)}\\
The algorithm descends quickly and converges in 20 steps, reaching $f(x) \approx 0.013$, very close to the true minimum of 0, and match the table in the lecture.
 
\begin{verbatim}
Time = 1  | x = 10.0000 | f(x) = 100.0000
Time = 2  | x =  8.0000 | f(x) =  64.0000
Time = 3  | x =  6.4000 | f(x) =  40.9600
...
Time = 18 | x =  0.2252 | f(x) =   0.0507
Time = 19 | x =  0.1801 | f(x) =   0.0325
Time = 20 | x =  0.1441 | f(x) =   0.0208

 Converged after 20 iterations, f(x) = 0.0133
\end{verbatim}
\vspace{0.5em}

\textbf{Small learning rate ($lr = 0.01$)}\\
The progress is steady but much slower. After 150 steps the threshold is met, but the final value is still $f(x) \approx 0.23$, farther from 0 than the $lr = 0.1$ case.
 
\begin{verbatim}
Time = 1   | x = 10.0000 | f(x) = 100.0000
Time = 2   | x =  9.8000 | f(x) =  96.0400
Time = 3   | x =  9.6040 | f(x) =  92.2368
...
Time = 148 | x =  0.5131 | f(x) =   0.2633
Time = 149 | x =  0.5029 | f(x) =   0.2529
Time = 150 | x =  0.4928 | f(x) =   0.2429
 
 Converged after 150 iterations, f(x) = 0.2333
\end{verbatim}
 
\vspace{0.5em}
\textbf{Very small learning rate ($lr = 0.001$)} \\
The steps are so small that it takes 923 iterations to converge, with $f(x) \approx 2.48$. The result is far from the actual minimum of 0, so this learning rate is too small to be useful.
 
\begin{verbatim}
Time = 1   | x = 10.0000 | f(x) = 100.0000
Time = 2   | x =  9.9800 | f(x) =  99.6004
Time = 3   | x =  9.9600 | f(x) =  99.2024
...
Time = 921 | x =  1.5853 | f(x) =   2.5130
Time = 922 | x =  1.5821 | f(x) =   2.5030
Time = 923 | x =  1.5789 | f(x) =   2.4930
 
 Converged after 923 iterations, f(x) = 2.4830
\end{verbatim}

\section*{Comparison}
Finally, I create a table of the effect of different learning rates on the gradient descent's performance. 

\begin{table}[h]
\centering
\begin{tabular}{@{}cccc@{}}
\toprule
\textbf{Learning Rate} & \textbf{Iterations} & \textbf{Final $f(x)$} & \textbf{Result} \\ \midrule
1.0   & 1   & 100.0000 & Diverged \\
0.1   & 20  & 0.0133   & Converged \\
0.01  & 150 & 0.2333   & Converged \\
0.001 & 923 & 2.4830   & Converged \\ \bottomrule
\end{tabular}
\caption{Comparison of gradient descent with different learning rates.}
\end{table}
 
\end{document}